\section{Scaling 之后：缺失能力、效率与 hallucination}

应用扩张并没有消除基础问题。课程把“未来”拆成两类：可以继续工程化的能力，如 on-device efficiency、personalization 与 monitoring；以及定义仍不清晰的问题，如 world understanding、hallucination 和长期适应。读本章时要区分 scaling law 的经验趋势、产品愿望和可检验机制。

\subsection{下一步与缺口：更大模型只是一个方向}

slide 121 列出可见应用增长：foundation model 组合、agent、机器人、个性化、医疗与交互娱乐。slide 122 则列出缺失条件：reliable reasoning、causal understanding、continual learning、alignment 与效率。前者回答“还能部署到哪里”，后者回答“要可靠部署还缺什么”。

\lecturefigure{slide-121.jpg}{未来应用：agent、机器人、个性化、医疗与交互系统}{CS25 V6 official deck, p.~121}

\lecturefigure{slide-122.jpg}{未来缺口：推理、因果、多模态、持续学习与 alignment}{CS25 V6 official deck, p.~122}

\teachervoice{00:52:47--00:53:34，老师先列潜在应用，再迅速转向 missing ingredients。课堂提示：应用列表是方向，缺口列表才定义研究问题与部署风险。}

\begin{knowledgebox}{把 future slide 变成工程问题}
“个性化”需要 identity/memory policy；“robotics”需要 perception-action loop 与安全约束；“healthcare”需要 calibration、provenance 与临床验证；“continual learning”需要可更新状态和遗忘控制。只有写出输入、更新和 metric，未来愿景才从口号变成研究计划。
\end{knowledgebox}

\subsection{Minified/on-device：模型要进入受限硬件}

On-device LLM 受 memory、功耗、latency 与更新渠道约束。压缩可分为 quantization、pruning、distillation、low-rank adaptation 与 architecture redesign。若参数量为 $N$、每参数位宽为 $b$，仅权重存储约为
\[
M_{\mathrm{weight}}\approx \frac{Nb}{8}\ \text{bytes}.
\]
真实峰值还包括 activation、KV cache、runtime buffer 和 tokenizer/embedding。仅把权重量化到 4-bit 不代表整个系统缩小四倍。

\lecturefigure{slide-123.jpg}{Minified and on-device LLM：效率、隐私与可用性目标}{CS25 V6 official deck, p.~123}

\begin{importantbox}{端侧优化的 Pareto frontier}
需要同时比较质量、首 token latency、每 token latency、峰值内存、能耗、热设计和可更新性。一个 benchmark 分数相近但持续发热、无法加载长 context 的模型，并不构成可用的端侧方案。
\end{importantbox}

\subsection{Scaling 的递减收益与下一批轴}

经验 scaling law 常写成
\[
L(N,D,C)\approx L_{\infty}+aN^{-\alpha}+bD^{-\beta}+cC^{-\gamma},
\]
$N$ 是参数、$D$ 是数据、$C$ 是 compute，指数描述特定 regime 下的边际改善。该关系预测 average loss，不保证事实可靠性、工具使用或 long-horizon reasoning 同步改善。模型越大，训练与 serving 成本越高，数据质量和评估污染也更难控制。

\lecturefigure{slide-124.jpg}{Limits of Scaling：递减收益、数据瓶颈与 compute 成本}{CS25 V6 official deck, p.~124}

\lecturefigure{slide-125.jpg}{Scaling 之后：新架构、学习程序与理论方向}{CS25 V6 official deck, p.~125}

\teachervoice{00:53:34--00:55:23，老师指出单纯 scaling 已出现 diminishing returns，并把 curriculum、reasoning procedure、理论解释和替代架构列为下一批轴。讲义提醒：这不是“scale 已失效”，而是边际成本上升。}

\begin{warningbox}{Loss scaling 与 capability scaling 不同}
更低 cross-entropy 可能改善平均生成，却不保证某个稀有技能、事实可靠性或 agent safety 单调提升。Capability 往往依赖数据覆盖、elicitation、post-training 和评估阈值；不能从一条 loss curve 推出所有能力的时间表。
\end{warningbox}

\subsection{Hallucination basics：先不要把所有错误放进一个桶}

课程先列出常见含义：生成不真实内容、与来源冲突、缺乏支持、误导或错误行动。问题在于不同任务使用不同 reference：摘要有 source document，开放问答有现实事实，agent 还有 observation 与 action consequence。若定义不指定 reference 与模型可见信息，就无法比较 benchmark，也无法区分知识缺失、观察缺失和 planning failure。

\lecturefigure{slide-127.jpg}{Hallucination basics：错误信息、缺乏支持与误导输出}{CS25 V6 official deck, p.~127}

\lecturefigure{slide-128.jpg}{Hallucination 为什么重要：可靠性、专业领域与信任}{CS25 V6 official deck, p.~128}

\teachervoice{00:55:23--00:58:11，老师用 summarization、QA 与 agent 说明同一输出在不同 setting 中可能有不同标签。课堂提示：先定义 reference，再讨论 mitigation。}

\subsection{没有统一定义：world-model view}

课程引用的统一框架把 hallucination 解释为错误 world modeling，而不是所有错误行为。设 reference world 为 $W$，模型可见 world 为 $V$，冲突处理 policy 为 $P$。模型输出的命题集合为 $B_{\theta}(V,P)$，可将 hallucination 条件写成
\[
\exists b\in B_{\theta}(V,P):\quad b\not\models W,
\]
即模型形成或表达了与 reference world 不相容的 belief。若 belief 正确但 action plan 糟糕，则属于 planning error，不应自动叫 hallucination。

\lecturefigure{slide-130.jpg}{No Unified Definition：现有 hallucination 定义覆盖范围不同}{CS25 V6 official deck, p.~130}

\lecturefigure{slide-131.jpg}{核心洞见：hallucination 是 world-model error，而非所有输出错误}{CS25 V6 official deck, p.~131}

\lecturefigure{slide-132.jpg}{统一框架：reference world $W$、visible world $V$ 与 policy $P$}{CS25 V6 official deck, p.~132}

\teachervoice{00:58:11--01:00:22，老师把定义依赖压缩为 $W,V,P$：什么是真实、模型看到了什么、冲突如何解决。讲义提醒：这套 formalization 把 assumption 暴露出来，但不会自动解决 truth source 的争议。}

\begin{knowledgebox}{首次使用：world model 在这里指什么}
这里的 world model 不是专指视频生成器，而是系统对相关世界状态、关系与规律的内部 belief/representation。Hallucination 被限定为该表示与 reference world 冲突。若模型根本未观察到信息，问题可能是 observability；若观察正确但策略失败，问题可能是 planning。
\end{knowledgebox}

\subsection{跨 setting 统一，但不要抹平误差类型}

slide 133 将 summarization、QA、agent 等 setting 映射到同一框架；slide 134 强调 not all errors are hallucinations。统一价值在于明确 reference、visible evidence 和 conflict rule，而不是把 benchmark 强行合并。不同 setting 的损失函数、可观测性和容错标准仍不同。

\lecturefigure{slide-133.jpg}{Unifying settings：摘要、QA 与 agent 的 $W,V,P$ 映射}{CS25 V6 official deck, p.~133}

\lecturefigure{slide-134.jpg}{Not all errors are hallucinations：planning、reward 与 belief error 的区分}{CS25 V6 official deck, p.~134}

\begin{warningbox}{同一个错误标签可能隐藏三条根因链}
\textbf{Knowledge error}：参数中没有正确规律；\textbf{Observation error}：正确证据未进入 context 或 retrieval 失败；\textbf{Decision error}：belief 正确但规划/动作错误。三者需要不同 mitigation：训练数据、检索接口和 planner/verifier 不能互相替代。
\end{warningbox}

\subsection{为什么定义重要：评估、干预与 synthetic environment}

明确 $W,V,P$ 后，可以构造 controlled environment：固定真值、改变可见证据、指定冲突规则，系统性生成 hallucination case。这样可比较 RAG、calibration、abstention、self-check 与 tool verification。但 synthetic benchmark 只覆盖建模的环境，仍可能遗漏现实中的开放世界、权限和时效性。

\lecturefigure{slide-135.jpg}{定义的作用一：暴露假设并比较不同方法}{CS25 V6 official deck, p.~135}

\lecturefigure{slide-136.jpg}{定义的作用二：构造可控环境与大规模合成案例}{CS25 V6 official deck, p.~136}

\lecturefigure{slide-137.jpg}{Hallucination big picture：world modeling、setting 与 evaluation}{CS25 V6 official deck, p.~137}

\teachervoice{01:00:22--01:03:52，老师强调 formal definition 的工程价值：可以比较 setting、生成 synthetic cases、训练和评估 mitigation；同时所有结论都依赖真值与可见世界的建模方式。}

\begin{importantbox}{Hallucination mitigation 的诊断顺序}
先确认 reference world 是否可靠，再检查证据是否进入 $V$，然后判断模型 belief 是否错误，最后检查 planner/action。若一开始就让模型“再想一遍”，可能只是在错误 evidence 上生成更长解释。
\end{importantbox}

例如一个医疗 agent 推荐错误药物，至少有四条不同根因：知识库本身过期（$W$ reference 错）；检索未取回过敏史（$V$ 不完整）；模型把已见证据解释错（world-model belief 错）；belief 正确但 action policy 忽略禁忌（planning/alignment error）。四种情况的修复分别是更新 source、改善 retrieval/permissions、训练或校准模型、修改 planner 与 safety rule。把它们都标成 hallucination 会让 root-cause analysis 失效，也会使 benchmark 看似统一、实际不可行动。

\subsection{本章小结}

Scaling 仍有效，但不自动解决效率、长期适应和可靠性。端侧部署要求完整资源核算，scaling law 只描述特定 loss regime。Hallucination 则必须先明确 $W,V,P$，把 belief error 与 observation/planning error 分开；精确定义是诊断和评估的入口，不是答案本身。

